% !PRINTED 86
\weberpage{86}
%<<CONT
\[
\frac{e_{1,h}i_{1,h}x_1}{i_1} + \frac{e_{2,h}i_{2,h}x_2}{i_2} + \cdots \frac{e_{\nu,h}i_{\nu,h}x_\nu}{i_\nu}
\]
ganze Zahlen sein müssen, oder, mit Rücksicht auf die Bezeichnung (2), mit den Congruenzen:
\weq{17}{e_{1,h}i_{1,h}i'_1x_1 + e_{2,h}i_{2,h}i'_2x_2 + \cdots e_{\nu,h}i_{\nu,h}i'_\nu x_\nu \equiv 0\ (\mathrm{mod}\ J).}
Diese Congruenzen sind, wie man sieht, immer erfüllt, wenn
$x_1$ durch $i_1$, $x_2$ durch $i_2 \dots$, $x_\nu$ durch $i_\nu$ theilbar ist, und die
Forderung 2. reducirt sich also jetzt darauf:

    5. \gesperrt{dass die $e_{k,h}$ so zu bestimmen sind, dass die Congruenzen (17) für keine anderen als den Congruenzen $x_k \equiv 0\ (\mathrm{mod}\ i_k)$ genügende Werthe der $x$ befriedigt sein sollen.}

    Dazu kommt noch als Umformung der Forderung 3.:

    6. Die Zahlen
\[
c_{1,h}e_{1,h},\ c_{2,h}e_{2,h},\ \dots\ c_{\nu,h}e_{\nu,h}
\]
\gesperrt{dürfen nicht alle durch $q_h$, wenn $\varkappa > 1$, und
durch $(q_h - 1)$, wenn $\varkappa = 1$ ist, theilbar sein.}

    Die Forderung 6. enthält wegen 4. keine Unmöglichkeit.

    \gesperrt{Durch 5. und 6. sind die nothwendigen und hinreichenden Bedingungen für die Zahlen $e_{h,k}$ ausgedrückt.}

    Es handelt sich also jetzt noch um die Frage, ob und unter welchen Voraussetzungen diese Forderungen durch geeignete Wahl der $e_{h,k}$ befriedigt werden können.
% [sic: Weber switches in these two sentences from the defined e_{k,h} to e_{h,k}; e_{k,h} is intended.]

    Diese Frage wird dadurch sehr vereinfacht, dass sich die
Congruenz (17) in mehrere spalten lässt. Die Invarianten $i_k$
sind, wie wir angenommen haben, Primzahlpotenzen. Es möge
eine von diesen Primzahlen mit $p$ bezeichnet sein, und es sei
\weq{18}{i_1 = p^{\pi_1},\ i_2 = p^{\pi_2},\ \dots\ i_\varrho = p^{\pi_\varrho},}
während die übrigen $i_k$, wenn noch solche vorhanden sind, nicht
mehr durch $p$ theilbar sein sollen. Den grössten der Exponenten $\pi_1, \pi_2, \dots \pi_\varrho$ wollen wir mit $\pi$ bezeichnen.

    Dann ist
\[
J = p^\pi J',
\]
und $J'$ ist nicht mehr durch $p$ theilbar. Es ist ferner
\[
i'_1 = p^{\pi-\pi_1}J',\ i'_2 = p^{\pi-\pi_2}J',\ \dots\ i'_\varrho = p^{\pi-\pi_\varrho}J',
\]
während $i'_{\varrho+1}, \dots i'_\nu$ durch $p^\pi$ theilbar sind.
